% Vietnamese translation for OI Public Library
% Source-PDF: IOI中国国家候选队论文/2007/day1/9.王欣上《浅谈基于分层思想的网络流算法》.pdf
\documentclass[11pt]{article}
\usepackage{oipl-vn}

\usepackage{array}

\newcommand{\Oh}{\mathcal{O}}
\newcommand{\level}{\operatorname{level}}
\newcommand{\throughput}{\operatorname{throughput}}
\newcommand{\outdeg}{\operatorname{outdeg}}
\newcommand{\inftycap}{\mathrm{INF}}

\title{Bàn sơ về các thuật toán luồng mạng dựa trên tư tưởng phân tầng}
\author{Wang Xinshang\\Trường Trung học Diên An, Thượng Hải}
\date{2007}

\begin{document}
\maketitle

\origtitle{Bàn sơ về các thuật toán luồng mạng dựa trên tư tưởng phân tầng}
\sourcepath{IOI China candidate team papers / 2007 / day 1 / Wang Xinshang}

\begin{abstract}
Bài viết giới thiệu tương đối chi tiết ba thuật toán luồng mạng dựa trên khái
niệm đồ thị phân tầng, đồng thời dùng các bài toán mẫu để minh họa ưu thế của
thuật toán Dinic trong các kỳ thi tin học.
\end{abstract}

\paragraph{Từ khóa.}
Đồ thị phân tầng; luồng mạng; thuật toán cơ bản; ứng dụng; MPLA; Dinic; MPM.

\section{Dẫn nhập}

Lý thuyết đồ thị là một ngành vừa cổ điển vừa trẻ trung, và chiếm tỉ trọng rất
lớn trong các kỳ thi tin học. Trong đó, các thuật toán luồng mạng xuất hiện
khá thường xuyên. Kiến thức về luồng mạng rất phong phú: từ định lý luồng cực
đại--cắt cực tiểu cơ bản, các biến thể của mô hình mạng, cho tới nhiều kỹ
thuật tối ưu cho thuật toán đẩy luồng trước. Muốn ứng phó với các dạng bài
khác nhau, học sinh cần tích lũy dần những kiến thức này.

Khi các kỳ thi phát triển, độ khó và phạm vi kiểm tra cũng tăng lên. Với nhiều
bài toán mới, chỉ nắm thuật toán tăng luồng thô sơ nhất là chưa đủ. Bài viết
này trình bày ba thuật toán luồng mạng dựa trên tư tưởng phân tầng, hướng tới
các bài toán có dữ liệu lớn hơn, và thông qua so sánh để làm rõ cách dùng
chúng khi giải bài. Những kiến thức nền như định lý luồng cực đại--cắt cực
tiểu sẽ không được chứng minh lại ở đây; độc giả có thể tìm thấy trong nhiều
tài liệu khác về luồng mạng.

\section{Các khái niệm chuẩn bị}

\subsection{Đồ thị dư}

Cho một mạng luồng \(G_1=(V_1,E_1)\), nguồn \(s\), đích \(t\), hàm sức chứa
\(c\), và một hàm luồng khả thi \(f\) trên mạng đó. Đồ thị dư tương ứng
\(G_2=(V_2,E_2)\) được định nghĩa như sau. Tập đỉnh không đổi:
\[
  V_2=V_1.
\]
Với mỗi cạnh có hướng \((u,v)\in E_1\):
\begin{itemize}
  \item nếu \(f(u,v)<c(u,v)\), đưa cạnh \((u,v)\) vào \(E_2\) với trọng số
  \[
    g(u,v)=c(u,v)-f(u,v);
  \]
  \item nếu \(f(u,v)>0\), đưa cạnh ngược \((v,u)\) vào \(E_2\) với trọng số
  \[
    g(v,u)=f(u,v).
  \]
\end{itemize}

Vì vậy, mỗi cạnh của mạng ban đầu biến thành một hoặc hai cạnh trong đồ thị
dư. Mỗi cạnh dư cho biết ta còn có thể tăng luồng theo hướng ấy trong mạng
gốc; trọng số \(g(a,b)\) chính là lượng luồng lớn nhất có thể tăng theo hướng
từ \(a\) tới \(b\) ở thời điểm hiện tại. Do đó, mọi đường đi đơn từ \(s\) tới
\(t\) trong đồ thị dư tương ứng với một đường tăng luồng, và giá trị nhỏ nhất
của các trọng số trên đường đi là lượng luồng lớn nhất có thể tăng trong một
lần theo đường ấy.

\subsection{Mức của đỉnh}

Trong đồ thị dư, độ dài đường đi ngắn nhất từ nguồn \(s\) tới đỉnh \(u\) được
gọi là \term{mức} của \(u\), ký hiệu \(\level(u)\). Nguồn có mức \(0\). Nếu
một đỉnh không đi tới được từ nguồn trong đồ thị dư, mức của nó không được
xác định trong pha hiện tại.

Một cách nhìn đơn giản là chia các đỉnh thành từng lớp:
\[
  L_i=\{u\mid \level(u)=i\},\qquad i=0,1,2,\ldots .
\]
Mỗi lần chạy BFS từ nguồn trên đồ thị dư, ta thu được các lớp này.

\subsection{Đồ thị phân tầng}

Đồ thị phân tầng \(G_3=(V_3,E_3)\) được dựng trên đồ thị dư \(G_2=(V_2,E_2)\).
Với một cạnh dư \((u,v)\), ta giữ cạnh ấy trong \(E_3\) khi và chỉ khi
\[
  \level(u)+1=\level(v).
\]
Tập đỉnh \(V_3\) gồm các đỉnh kề với ít nhất một cạnh trong \(E_3\), cùng với
nguồn và đích khi cần xét đường tăng luồng.

Trực giác là: đồ thị phân tầng là ``đồ thị các đường ngắn nhất'' xây trên đồ
thị dư. Nếu bắt đầu từ nguồn và đi theo các cạnh trong đồ thị phân tầng, thì
sau \(k\) cạnh ta luôn đứng ở một đỉnh có khoảng cách ngắn nhất bằng \(k\) so
với nguồn trong đồ thị dư.

\subsection{Luồng chặn}

Xét một luồng khả thi \(f\) trong mạng. Nếu trong đồ thị phân tầng của đồ thị
dư tương ứng không còn đường tăng luồng từ \(s\) tới \(t\), ta gọi \(f\) là
\term{luồng chặn} của đồ thị phân tầng ấy.

Nói cách khác, luồng chặn không nhất thiết là luồng cực đại của toàn mạng, mà
chỉ chặn hết mọi đường tăng luồng ngắn nhất trong pha hiện tại.

\section{Thuật toán tăng theo đường ngắn nhất MPLA}

MPLA là viết tắt của \emph{maximum path length augmenting} trong bản dịch này:
ta luôn làm việc trên các đường tăng luồng ngắn nhất bằng cách dựng đồ thị
phân tầng. Nội dung chính của thuật toán là lặp lại việc dựng tầng và tăng
luồng trong đồ thị tầng cho tới khi không còn đường tăng luồng.

\subsection{Các bước thuật toán}

\begin{enumerate}
  \item Khởi tạo luồng, từ đó có đồ thị dư ban đầu.
  \item Chạy BFS trên đồ thị dư để tính mức các đỉnh, tức dựng đồ thị phân
  tầng. Nếu đích không xuất hiện trong các mức, thuật toán kết thúc.
  \item Trong đồ thị phân tầng, liên tục dùng BFS tìm đường tăng luồng và tăng
  luồng theo đường đó cho tới khi đồ thị phân tầng không còn đường tăng luồng.
  \item Quay lại bước 2.
\end{enumerate}

Một lần thực hiện các bước 2 và 3 được gọi là một \term{pha}. Trong mỗi pha,
ta dựng đồ thị phân tầng từ đồ thị dư hiện tại, sau đó tìm một luồng chặn
trong đồ thị phân tầng ấy. Sau khi không còn đường tăng luồng trong tầng,
thuật toán sang pha tiếp theo.

Khi cài đặt, không cần thật sự tạo một cấu trúc đồ thị mới cho \(G_3\). Chỉ cần
gán \(\level\) cho từng đỉnh, rồi khi duyệt cạnh \((u,v)\) kiểm tra điều kiện
\(\level(u)+1=\level(v)\).

\subsection{Số pha}

\paragraph{Định lý.}
Trong một mạng có \(n\) đỉnh, thuật toán MPLA có nhiều nhất \(n\) pha.

\paragraph{Chứng minh.}
Giả sử sau khi dựng đồ thị phân tầng, độ dài đường ngắn nhất từ \(s\) tới
\(t\) trong đồ thị dư là \(k\). Chia các đỉnh đi tới được từ nguồn thành
\(k+1\) lớp:
\[
  L_i=\{u\mid \level(u)=i\},\qquad 0\le i\le k.
\]
Trong đồ thị dư lúc đó, các cạnh có thể được nhìn theo hai loại quan trọng:
\begin{itemize}
  \item cạnh loại một đi từ \(L_i\) tới \(L_{i+1}\);
  \item cạnh loại hai đi từ \(L_i\) tới \(L_j\) với \(j\le i\).
\end{itemize}
Đồ thị phân tầng chỉ chứa cạnh loại một. Vì vậy mọi đường tăng luồng trong
tầng đều gồm đúng các cạnh loại một và có độ dài \(k\).

Mỗi lần tăng luồng theo một đường trong tầng, ít nhất một cạnh trên đường bị
bão hòa và bị xóa khỏi đồ thị dư theo hướng thuận. Đồng thời, các cạnh ngược
của đường tăng luồng được thêm vào đồ thị dư; những cạnh ngược này đều đi từ
lớp sâu hơn về lớp nông hơn, tức thuộc loại hai.

Khi đã tìm được luồng chặn, không còn đường nào chỉ đi qua các cạnh loại một
từ \(L_0\) tới \(L_k\). Mọi đường mới từ \(s\) tới \(t\), nếu còn tồn tại trong
đồ thị dư, bắt buộc phải dùng ít nhất một cạnh loại hai để lùi lên một lớp
trước khi đi tiếp xuống các lớp sau. Vì đã phải lùi lại, số cạnh loại một cần
đi sau đó tăng lên, nên tổng độ dài đường đi lớn hơn \(k\).

Do đó, độ dài đường tăng luồng ngắn nhất tăng nghiêm ngặt sau mỗi pha. Độ dài
nhỏ nhất có thể là \(1\), lớn nhất có thể là \(n-1\), cộng thêm lần cuối cùng
khi đích không còn đạt được từ nguồn, nên số pha không vượt quá \(n\).

\subsection{Phân tích độ phức tạp}

Tách chi phí thành hai phần: dựng đồ thị phân tầng và tìm đường tăng luồng
trong một tầng.

Mỗi pha dựng tầng bằng một BFS trên đồ thị dư, chi phí \(\Oh(m)\). Vì có nhiều
nhất \(n\) pha, tổng chi phí dựng tầng là
\[
  \Oh(nm).
\]

Trong một pha, mỗi lần tăng luồng làm bão hòa ít nhất một cạnh của đồ thị
phân tầng. Đồ thị phân tầng có nhiều nhất \(m\) cạnh, nên số lần tăng luồng
trong một pha không vượt quá \(m\). Với MPLA, mỗi lần lại dùng BFS để tìm
đường tăng luồng; chi phí tìm đường là \(\Oh(m+n)\), và chi phí sửa luồng dọc
đường là \(\Oh(n)\). Vì vậy chi phí mỗi pha là
\[
  \Oh\bigl(m(m+n)\bigr)=\Oh(m^2)
\]
trong các mạng thường xét với \(m\ge n\). Tổng chi phí tăng luồng qua mọi pha
là \(\Oh(nm^2)\), và đây cũng là độ phức tạp tổng quát của MPLA.

\section{Thuật toán Dinic}

\subsection{Các bước thuật toán}

Dinic cũng chia quá trình thành các pha trên đồ thị phân tầng. Khác biệt then
chốt so với MPLA là: trong mỗi pha, Dinic dùng một quá trình DFS để tìm luồng
chặn, thay vì nhiều lần BFS độc lập.

\begin{enumerate}
  \item Khởi tạo luồng và đồ thị dư.
  \item Chạy BFS trên đồ thị dư để tính mức các đỉnh. Nếu không tới được đích,
  thuật toán kết thúc.
  \item Dùng DFS trên đồ thị phân tầng để tăng luồng cho tới khi thu được
  luồng chặn.
  \item Quay lại bước 2.
\end{enumerate}

Quá trình DFS trong một pha có thể mô tả bằng một đường hiện tại \(p\), ban
đầu chỉ gồm nguồn \(s\). Ký hiệu \(p.\mathrm{top}\) là đỉnh cuối của đường.

\begin{verbatim}
p <- [s]
while outdeg(s) > 0 do
    u <- p.top
    if u != t then
        if outdeg(u) > 0 then
            choose a level edge (u, v)
            push v to the end of p
        else
            delete u from p and from the level graph
            delete all level-graph edges incident to u
    else
        augment along p and delete saturated edges on p
        move p.top back to the last reachable vertex on p
end while
\end{verbatim}

Nếu đỉnh cuối của \(p\) là đích, ta đã có một đường tăng luồng. Sau khi tăng
luồng, ít nhất một cạnh trên đường bị bão hòa; đường hiện tại được lùi về đỉnh
xa nhất trên \(p\) vẫn còn đi tới được từ nguồn theo các cạnh chưa bị xóa.

Nếu đỉnh cuối \(u\) không phải đích, có hai khả năng. Nếu còn cạnh tầng đi ra
từ \(u\), ta đi tiếp theo cạnh đó. Nếu không còn cạnh đi ra, \(u\) không thể
nằm trên bất kỳ đường tăng luồng nào của phần đồ thị tầng còn lại; khi đó xóa
\(u\), xóa các cạnh kề với \(u\), rồi lùi đường hiện tại một đỉnh.

Quá trình trên kết thúc khi mọi cạnh đi ra từ nguồn trong đồ thị tầng đã bị
xóa. Khi đó đồ thị tầng không còn đường từ \(s\) tới \(t\), tức ta đã tìm được
một luồng chặn.

\subsection{Minh họa bằng lời}

Trong một đồ thị phân tầng, Dinic luôn giữ một đường hiện tại từ nguồn đi dần
về phía đích:
\[
  s \longrightarrow \cdots \longrightarrow u.
\]
Khi còn cạnh hợp lệ từ \(u\), đường được kéo dài. Khi gặp đích, ta đẩy thêm
luồng trên cả đường; các cạnh bão hòa bị loại khỏi pha hiện tại. Khi gặp một
đỉnh cụt, đỉnh đó bị xóa vì không thể đóng góp cho đường tăng luồng nào nữa.

Do chỉ đi theo cạnh từ lớp \(i\) sang lớp \(i+1\), DFS không tạo chu trình
trong một pha. Việc xóa cạnh và xóa đỉnh là nguyên nhân chính giúp Dinic tránh
lặp lại nhiều phép tìm đường tốn kém như MPLA.

\subsection{Phân tích độ phức tạp}

Lập luận về số pha giống MPLA: sau mỗi pha, độ dài đường tăng luồng ngắn nhất
tăng nghiêm ngặt, nên Dinic có nhiều nhất \(n\) pha. Tổng chi phí dựng đồ thị
phân tầng vẫn là \(\Oh(nm)\).

Xét chi phí DFS trong một pha.

\paragraph{Sửa luồng và lùi sau khi tăng.}
Trong một pha, số lần tăng luồng không vượt quá \(m\), vì mỗi lần tăng làm bão
hòa ít nhất một cạnh tầng. Mỗi lần sửa luồng dọc đường mất \(\Oh(n)\), và việc
lùi đường hiện tại cũng mất \(\Oh(n)\). Vì vậy phần này tốn
\[
  \Oh(mn)
\]
cho một pha.

\paragraph{Đi tới và lùi khi duyệt DFS.}
Mỗi lần DFS đi tới, nó đi theo một cạnh từ lớp hiện tại sang lớp kế tiếp. Một
đường đơn trong tầng có độ dài nhiều nhất \(n\). Khi gặp đỉnh cụt, thuật toán
xóa đỉnh đó; số lần lùi kiểu này trong một pha nhiều nhất \(n\).

Trong trường hợp xấu nhất, có \(n\) lần đi tới bị xóa bỏ cùng với các đỉnh cụt.
Các lần đi tới còn lại đều góp phần vào những đường tăng luồng đã tìm được.
Có nhiều nhất \(m\) đường tăng luồng trong một pha, mỗi đường dài nhiều nhất
\(n\), nên tổng số bước đi tới là \(\Oh(mn)\). Cộng thêm các lần lùi do đỉnh
cụt, chi phí duyệt DFS trong một pha vẫn là \(\Oh(mn)\).

Vì thế, chi phí tìm luồng chặn trong một pha là \(\Oh(mn)\). Nhân với nhiều
nhất \(n\) pha, độ phức tạp tổng quát của Dinic là
\[
  \Oh(mn^2).
\]

Trong thực tế thi đấu, Dinic thường nhanh hơn đánh giá tổng quát này rất nhiều,
đặc biệt trên các mạng hai phía và các mạng có sức chứa nguyên nhỏ.

\section{Ứng dụng của Dinic trong thi tin học}

\subsection{Bài toán 1: Lợi nhuận lớn nhất}

\paragraph{Mô tả bài toán.}
Có \(N\) trạm trung chuyển tín hiệu liên lạc. Xây trạm \(i\) tốn chi phí
\(P_i\). Ngoài ra có \(M\) nhóm người dùng. Nhóm \(j\) sẽ dùng hai trạm
\(A_j\) và \(B_j\) để liên lạc; nếu phục vụ nhóm này, công ty thu được
\(C_j\). Hãy chọn xây một số trạm sao cho lợi nhuận ròng là lớn nhất.

Giới hạn trong bài gốc: \(N\le 5000\), \(M\le 50000\).

\paragraph{Mô hình hóa ban đầu.}
Biểu diễn trạm \(i\) bằng một đỉnh có trọng số \(-P_i\). Biểu diễn nhóm người
dùng \(j\) bằng một cạnh nối hai đỉnh \(A_j,B_j\), có trọng số \(C_j\). Ta cần
chọn một tập đỉnh sao cho tổng trọng số các đỉnh đã chọn cộng với tổng trọng
số các cạnh có cả hai đầu mút trong tập là lớn nhất.

Hãy tưởng tượng ban đầu ta định xây mọi trạm và phục vụ mọi nhóm người dùng:
ta thu \(\sum C_j\) và trả \(\sum P_i\). Mỗi cạnh lợi nhuận có thể dùng một
phần giá trị của nó để bù chi phí cho hai đầu mút. Nếu một đỉnh vẫn âm sau khi
mọi cạnh kề đã góp hết, giữ nó là không có lợi; ta nên bỏ trạm đó và bỏ các
nhóm người dùng cần trạm ấy. Việc bỏ một đỉnh có thể khiến các đỉnh khác mất
phần bù chi phí và cũng trở nên không có lợi. Quá trình này chính là một bài
toán cắt tối thiểu.

\paragraph{Mạng luồng.}
Dựng mạng hai phía như sau.
\begin{itemize}
  \item Với mỗi nhóm người dùng \(j\), tạo đỉnh \(X_j\) và nối
  \(S\to X_j\) với sức chứa \(C_j\).
  \item Với mỗi trạm \(i\), tạo đỉnh \(Y_i\) và nối \(Y_i\to T\) với sức chứa
  \(P_i\).
  \item Với nhóm \(j\), nối \(X_j\to Y_{A_j}\) và \(X_j\to Y_{B_j}\), mỗi cạnh
  có sức chứa \(\inftycap\).
\end{itemize}

Một lát cắt hữu hạn không thể cắt các cạnh \(\inftycap\), do đó nếu giữ lợi
nhuận của nhóm \(j\) ở phía nguồn thì cũng phải giữ cả hai trạm tương ứng ở
phía nguồn. Giá trị lát cắt tối thiểu bằng
\[
  \text{chi phí các trạm được giữ}
  + \text{lợi nhuận ban đầu của các nhóm bị bỏ}.
\]
Vì tổng lợi nhuận ban đầu \(\sum C_j\) là hằng số, đáp án là
\[
  \sum_{j=1}^{M} C_j - \operatorname{maxflow}.
\]

\paragraph{So sánh thuật toán.}
Với dữ liệu lớn, MPLA thường không qua được các bộ kiểm thử lớn. Tác giả ghi
nhận thời gian đại diện như sau:
\[
\begin{array}{c|ccc}
\text{Thuật toán} & \text{Test 1--8} & \text{Test 9} & \text{Test 10}\\
\hline
\text{MPLA} & <0.1\text{s} & >30\text{s} & >30\text{s}\\
\text{Dinic} & <0.03\text{s} & 0.40\text{s} & 0.37\text{s}\\
\text{Đẩy luồng trước} & <0.03\text{s} & 0.53\text{s} & 0.51\text{s}
\end{array}
\]

Một hướng khác là dựng luồng ban đầu bằng tham lam, chẳng hạn ưu tiên các đỉnh
bậc nhỏ trong đồ thị hai phía. Tuy nhiên cách này phụ thuộc nhiều vào đặc điểm
đồ thị của bài. Dinic ngắn gọn hơn, ít sai hơn, và trong bài này đủ nhanh để
không cần dựa vào luồng ban đầu đặc biệt.

\paragraph{Nhận xét.}
Trên mạng hai phía của bài này, Dinic hơi nhanh hơn đẩy luồng trước theo số
liệu của tác giả. Điều đó không có nghĩa Dinic luôn nhanh hơn mọi biến thể
đẩy luồng trước; mỗi thuật toán có loại đồ thị thuận lợi riêng. Lợi thế lớn
của Dinic trong thi đấu là cân bằng tốt giữa tốc độ, độ dễ hiểu và độ dễ cài.

\subsection{Bài toán 2: Trò chơi ma trận}

\paragraph{Mô tả bài toán.}
Cho các số \(R_1,\ldots,R_n\) và \(C_1,\ldots,C_m\). Hỏi có tồn tại ma trận
\(0\)-\(1\) kích thước \(n\times m\) sao cho hàng \(i\) có đúng \(R_i\) ô bằng
\(1\), cột \(j\) có đúng \(C_j\) ô bằng \(1\). Ngoài ra, ở mỗi hàng và mỗi
cột có nhiều nhất một ô bị chỉ định bắt buộc bằng \(0\). Giới hạn:
\(n,m\le 1000\), mỗi bộ dữ liệu có nhiều nhất \(10\) test.

\paragraph{Mô hình luồng trực tiếp.}
Dựng mạng hai phía:
\begin{itemize}
  \item tạo đỉnh \(X_i\) cho hàng \(i\), nối \(S\to X_i\) với sức chứa \(R_i\);
  \item tạo đỉnh \(Y_j\) cho cột \(j\), nối \(Y_j\to T\) với sức chứa \(C_j\);
  \item nếu ô \((i,j)\) không bị bắt buộc bằng \(0\), nối \(X_i\to Y_j\) với
  sức chứa \(1\).
\end{itemize}
Nếu cạnh \(X_i\to Y_j\) mang một đơn vị luồng, đặt ô \((i,j)=1\). Bài toán có
lời giải khi và chỉ khi luồng cực đại bằng tổng số ô bằng \(1\), tức bằng
\(\sum_i R_i=\sum_j C_j\).

Cách dựng này có \(\Oh(nm)\) cạnh. Với MPLA chỉ đạt được một phần điểm; với
Dinic có thể đạt điểm cao hơn, nhưng vẫn chưa phải lời giải chuẩn của bài.

\paragraph{Một kiểm tra tham lam hỗ trợ.}
Tạm bỏ qua các ô bị ép bằng \(0\). Vì hoán vị toàn bộ các hàng hoặc các cột
không làm thay đổi bản chất bài toán, có thể sắp \(R_i\) và \(C_j\) giảm dần.
Khi xét lần lượt các cột, ta tưởng tượng mỗi hàng có thể cung cấp một ô \(1\)
cho các cột đầu tiên còn cần. Nếu một cột cần nhiều ô \(1\) hơn số hàng đang
có thể cung cấp, ta dùng các ô đã ``dư'' từ những cột trước. Nếu tới một cột
mà kể cả dùng hết phần dư vẫn không đủ, thì chắc chắn không tồn tại ma trận.

Kiểm tra này loại được nhiều trường hợp vô nghiệm trước khi chạy luồng. Trong
bài gốc, thêm kiểm tra đơn giản ấy rồi dùng Dinic có thể qua toàn bộ dữ liệu,
trong khi MPLA vẫn không đủ nhanh.

\paragraph{Nhận xét.}
Lời giải chuẩn của bài là một thuật toán tham lam đầy đủ. Dù vậy, ví dụ này
cho thấy trong thi đấu, một thuật toán luồng đủ mạnh có thể giúp ta đạt kết
quả tốt khi chưa kịp khai thác hết cấu trúc riêng của bài.

\section{Thuật toán MPM}

\subsection{Các bước thuật toán}

MPM cũng tìm luồng chặn theo từng pha trên đồ thị phân tầng, nhưng không dùng
DFS kiểu Dinic. Trước hết, định nghĩa \term{thông lượng} của một đỉnh
\(u\ne s,t\) trong đồ thị phân tầng:
\[
  \throughput(u)=
  \min\left(
    \sum_{(v,u)} g(v,u),
    \sum_{(u,w)} g(u,w)
  \right).
\]
Đây là lượng luồng lớn nhất có thể đi xuyên qua \(u\) nếu chỉ xét cân bằng
giữa tổng sức chứa đi vào và tổng sức chứa đi ra của \(u\).

Trong mỗi pha, MPM tìm luồng chặn như sau.
\begin{enumerate}
  \item Tìm một đỉnh \(u\) có \(\throughput(u)\) nhỏ nhất trong đồ thị phân
  tầng, không tính nguồn và đích.
  \item Từ \(u\), đẩy \(\throughput(u)\) đơn vị luồng về phía đích trong đồ
  thị phân tầng.
  \item Từ \(u\), kéo \(\throughput(u)\) đơn vị luồng từ phía nguồn trong đồ
  thị phân tầng.
  \item Xóa \(u\). Nếu đồ thị phân tầng chỉ còn nguồn và đích, hoặc không còn
  đường từ nguồn tới đích, pha kết thúc; nếu không, quay lại bước 1.
\end{enumerate}

Vì \(u\) là đỉnh có thông lượng nhỏ nhất, quá trình đẩy và kéo lượng
\(\throughput(u)\) sẽ không bị kẹt ở các đỉnh khác. Khi đẩy hoặc kéo, ta luôn
cố gắng làm bão hòa các cạnh đã dùng; tại mỗi bước chỉ có thể còn nhiều nhất
một cạnh chưa bão hòa hoàn toàn theo hướng đang xử lý.

\subsection{Phân tích độ phức tạp}

Trong một pha, mỗi vòng của MPM xóa một đỉnh, nên có nhiều nhất \(n\) vòng.
Nếu chọn đỉnh có thông lượng nhỏ nhất bằng cách quét tuyến tính, mỗi vòng tốn
\(\Oh(n)\). Mỗi lần đẩy và kéo xử lý nhiều nhất \(\Oh(n)\) cạnh chưa bão hòa
hoàn toàn, còn trong cả pha tổng số cạnh bị xóa là \(\Oh(m)\). Do đó chi phí
tìm luồng chặn trong một pha là
\[
  \Oh(n(n+n)+m)=\Oh(n^2+m).
\]
Trong cách đánh giá của bài gốc, lấy \(\Oh(n^2)\) cho thành phần chính của
việc tìm luồng chặn trên mỗi tầng. Kết hợp với nhiều nhất \(n\) pha, độ phức
tạp tổng quát của MPM được ghi là
\[
  \Oh(n^3).
\]

MPM có ý tưởng đẹp nhưng cài đặt rườm rà hơn Dinic. Với yêu cầu thi đấu thông
thường, lợi ích thực tế của nó không rõ bằng Dinic.

\section{Tổng kết}

Ba thuật toán có thể nhìn nhanh như sau:
\[
\begin{array}{c|c|c|c}
\text{Thuật toán} & \text{Ý tưởng trong mỗi pha} &
\text{Độ phức tạp} & \text{Đặc điểm}\\
\hline
\text{MPLA} & \text{nhiều BFS tìm đường tăng} &
\Oh(nm^2) & \text{dễ hiểu, hợp dữ liệu vừa và nhỏ}\\
\text{Dinic} & \text{DFS tìm luồng chặn} &
\Oh(mn^2) & \text{dễ cài, hợp dữ liệu lớn hơn}\\
\text{MPM} & \text{xử lý đỉnh có thông lượng nhỏ nhất} &
\Oh(n^3) & \text{ý tưởng rõ, cài đặt phức tạp hơn}
\end{array}
\]

So sánh tổng thể, Dinic là lựa chọn rất đáng ưu tiên trong các kỳ thi tin học:
nó dựa trên khái niệm phân tầng giống MPLA, nhưng tìm luồng chặn hiệu quả hơn
nhiều; đồng thời cách cài đặt vẫn đủ ngắn và ổn định để dùng trong điều kiện
thi.

\section*{Lời cảm ơn}

Tác giả cảm ơn thầy Yu Xiaoqing đã cung cấp đề bài và dữ liệu; cảm ơn Li
Tianyi, Hu Botao, Hu Gang, Shou Heming và huấn luyện viên Liu Rujia đã góp ý
cho bản thảo hoặc chương trình.

\section*{Tài liệu tham khảo}

\begin{enumerate}
  \item Thomas H. Cormen, Charles E. Leiserson, Ronald L. Rivest, Clifford
  Stein, \emph{Introduction to Algorithms}.
  \item M. H. Alsuwaiyel, \emph{Algorithms: Design Techniques and Analysis}.
  \item Liu Rujia, Huang Liang, \emph{The Art of Algorithms and Informatics
  Competitions}.
  \item Wang Xiaodong, \emph{Computer Algorithm Design Techniques and
  Analysis}.
  \item Fred Buckley, Marty Lewinter, \emph{A Friendly Introduction to Graph
  Theory}.
  \item Zhou Yuan, solution report for the 2006 Anhui provincial informatics
  selection contest.
\end{enumerate}

\section*{Phụ lục: đề bài mẫu}

\subsection*{Lợi nhuận lớn nhất}

Bài từ NOI 2006. Có \(N\) vị trí có thể xây trạm trung chuyển tín hiệu. Xây
trạm \(i\) tốn \(P_i\). Có \(M\) nhóm người dùng; nhóm \(i\) cần dùng hai trạm
\(A_i,B_i\) và đem lại lợi nhuận \(C_i\). Công ty được chọn xây một số trạm và
phục vụ một số nhóm người dùng. Hãy tối đa hóa
\[
  \text{lợi nhuận ròng}
  = \text{tổng lợi nhuận phục vụ} - \text{tổng chi phí xây trạm}.
\]

Dữ liệu: \(N\le 5000\), \(M\le 50000\), \(0\le C_i\le 100\),
\(0\le P_i\le 100\). Với \(80\%\) dữ liệu: \(N\le 200\), \(M\le 1000\).

\subsection*{Trò chơi ma trận}

Bài từ JSOI 2006. Với mỗi test, cho \(n,m\le 1000\), các tổng hàng
\(r_1,\ldots,r_n\), các tổng cột \(c_1,\ldots,c_m\), và một số ô bị bắt buộc
bằng \(0\). Mỗi hàng và mỗi cột có nhiều nhất một ô bị bắt buộc bằng \(0\).
Cần xác định có tồn tại ma trận \(0\)-\(1\) thỏa tất cả điều kiện hay không.

Với mỗi test, in \texttt{Yes} nếu tồn tại ma trận, ngược lại in \texttt{No}.
Các bộ dữ liệu đầu có \(1\le n,m\le 100\), các bộ cuối có
\(1\le n,m\le 1000\).

\end{document}
